%-------------------------------------------------------------------------------
% Technical appendices A--F, English translation of the source appendices.
%-------------------------------------------------------------------------------

\section{Knowledge-Graph Relation Types and Retrieval Admissibility}
\label{app:kg}

The Deterministic Context Retrieval procedure of \S\ref{sec:method-graphrag} builds context around a set of
core relations and explicitly distinguishes ``relations admissible into
retrieval'' from ``relations used only in post-processing''
as shown in Table~\ref{tab:kgrel}. The retrieval stage
obeys four constraints, GR-1 through GR-4: no inferred normative rule may enter
retrieval; only raw specification nodes may enter retrieval; no inferential
relation may enter retrieval; and all context must be traceable to a
specification source.

\begin{table}[tb]
\centering
\footnotesize
\setlength{\tabcolsep}{4pt}
\caption{Knowledge-graph relation types and retrieval admissibility.}
\label{tab:kgrel}
\begin{tabularx}{\columnwidth}{@{}l l X c@{}}
\toprule
Type & Relation & Role & In retrieval \\
\midrule
core & CONTAINS & containment hierarchy among specs, sections, rules & Yes \\
core & DEFINES & link from a term/field/definition passage to its defined
  object & Yes \\
core & REFERENCES & explicit reference to another spec/section/algorithm & Yes \\
core & APPLIES TO & link from a rule to the abstract concept it constrains
  & Yes \\
aux & AFFECTS & link from a rule to the concrete certificate field it affects;
  used by the field resolver & No \\
aux & DERIVED FROM & provenance link from a derived rule to its source section
  & No \\
inferred & OVERRIDES & a more specific rule overrides a general one (produced by
  the rule engine) & No \\
inferred & CONFLICTS WITH & potential conflict between rules (produced by the
  rule engine) & No \\
\bottomrule
\end{tabularx}
\end{table}

\section{Key Fields of the Structured IR}
\label{app:ir}

The system internally uses a JSON-serialized, machine-readable IR object.
Beyond the core four-tuple of \S\ref{sec:method-dual-layer}, its key fields are listed in Table~\ref{tab:irfields};
the remaining fields support provenance, normalization, matching, and downstream
product generation. The lintability decision in \S\ref{sec:lintability} directly
uses only the four fields obligation, assertion\_subject, enforcement\_phase, and
rule\_category.

\begin{table}[tb]
\centering
\footnotesize
\caption{Key fields of the structured IR.}
\label{tab:irfields}
\begin{tabularx}{\columnwidth}{@{}l l X@{}}
\toprule
Field & Role & Meaning \\
\midrule
subject & core & constrained certificate field or path \\
obligation & core & RFC~2119 deontic level, such as MUST/SHOULD/\dots \\
predicate & core & constraint relation, such as presence, equality,
  containment, or range \\
constraint & core & constraint value, pattern, or condition \\
rule\_category & key ext. & rule's semantic category, such as
  encoding\_constraint, definition, algorithm\_ref, \dots \\
assertion\_subject & key ext. & assertion subject, such as Certificate / CA /
  RelyingParty / external ecosystem \\
enforcement\_phase & key ext. & phase the constraint depends on
  such as encoding / runtime / external validation \\
source\_section & key ext. & specification source and section
  identifier \\
source\_span & key ext. & character/sentence span in the source text \\
evidence\_text & key ext. & source fragment supporting this IR record \\
context\_nodes & key ext. & definition, field metadata, or cited
  context nodes from Deterministic Context Retrieval \\
\bottomrule
\end{tabularx}
\end{table}

\section{Construction Details of the Restricted Code Space $\Tv$}
\label{app:tv}

\textbf{The seven components of the vocabulary $\mathcal{V}$.} $\mathcal{V}$ is
the disjoint union of seven finite sets
shown in Table~\ref{tab:vocab}: certificate fields $\mathcal{F}_{\mathrm{cert}}$,
DN fields $\mathcal{F}_{\mathrm{dn}}$ over the Subject and Issuer names, OID
constants $\mathcal{O}$, including zlint util.* extension and algorithm OIDs
together with well-known CABF and RFC OIDs, KeyUsage bits $\mathcal{B}_{\mathrm{KU}}$,
ExtKeyUsage bits $\mathcal{B}_{\mathrm{EKU}}$, ASN.1 encoding types
$\mathcal{E}_{\mathrm{ASN1}}$, and named regexes $\mathcal{R}_{\mathrm{regex}}$,
each an audited literal RE2 pattern. Table~\ref{tab:vocab} gives the size and
representative members of each component.
All components are finite sets; the LLM receives a full enumeration of
each component in the prompt, so any identifier outside $\mathcal{V}$ in its
output triggers an $\eta$ parse error.

\begin{table*}[tb]
\centering
\footnotesize
\caption{The seven frozen components of the restricted vocabulary $\mathcal{V}$,
with their sizes and representative members. The LLM may emit only identifiers
drawn from these closed sets; the full enumerations are released with the code
and data.}
\label{tab:vocab}
\begin{tabularx}{\textwidth}{@{}l l r X@{}}
\toprule
Symbol & Component & Size & Representative members \\
\midrule
$\mathcal{F}_{\mathrm{cert}}$ & Certificate fields & 52 &
  c.Version, c.NotAfter, c.KeyUsage, c.DNSNames, c.IPAddresses \\
$\mathcal{F}_{\mathrm{dn}}$ & DN fields over Subject / Issuer & 38 &
  Subject.CommonName, Subject.Organization, Subject.Country, Issuer.CommonName \\
$\mathcal{O}$ & OID constants & 99 &
  util.BasicConstOID, util.SubjectAlternateNameOID, util.AiaOID \\
$\mathcal{B}_{\mathrm{KU}}$ & KeyUsage bits & 9 &
  DigitalSignature, KeyEncipherment, CertSign, CRLSign \\
$\mathcal{B}_{\mathrm{EKU}}$ & ExtKeyUsage bits & 11 &
  ServerAuth, ClientAuth, CodeSigning, EmailProtection \\
$\mathcal{E}_{\mathrm{ASN1}}$ & ASN.1 encoding types & 11 &
  PrintableString, IA5String, UTF8String, BMPString \\
$\mathcal{R}_{\mathrm{regex}}$ & Named RE2 regexes & 28 &
  Re\_LDH\_Hostname, Re\_Rfc5321Mailbox, Re\_TorV3Onion \\
\bottomrule
\end{tabularx}
\end{table*}

\textbf{Representative atomic templates.} The atomic-template set $\mathcal{A}$
has $|\mathcal{A}|=151$ and is organized by semantic cluster, each with a typed
signature $\mathrm{sig}(a)$ described in \S\ref{subsec:tv}. Table~\ref{tab:atoms} excerpts
key atomic templates by cluster, and all 151 are released with the code and data.

\begin{table*}[tb]
\centering
\footnotesize
\caption{Representative atomic templates by semantic cluster.}
\label{tab:atoms}
\begin{tabularx}{\textwidth}{@{}l p{4.2cm} l X@{}}
\toprule
Cluster & Atomic template & Signature & Semantics \\
\midrule
I Extension presence & ExtPresent / ExtCritical &
  $(\mathcal{O})$ & extension OID present / present and critical bit set \\
II Macro attributes & IsCA / KeyUsageHas /
  ExtKeyUsageHas & $()$ / $(\mathcal{B}_{\mathrm{KU}})$ /
  $(\mathcal{B}_{\mathrm{EKU}})$ & BasicConstraints.cA$=$true / KeyUsage bit /
  EKU item \\
III Field value/form & FieldEq / FieldMatchesRegex /
  FieldEncodedAs & $(\mathcal{F}, \cdot)$ & field equals literal /
  matches named regex / ASN.1 encoding type \\
IV List/byte-level & ListAllMatch / OidListContains /
  BytesContainsOidDer & $(\mathcal{F}, \cdot)$ & list all satisfy
  subtree / contains named OID / bytes contain the OID's DER \\
V NameConstraints & SubtreeIPListAny\-HasOctetCount\-AndNotAllZero &
  $(\mathcal{F}, \mathbb{Z})$ & NC IP subtree has an entry of given octet count
  and not all-zero \\
VI Special structure & DomainComponentOrdered & $()$ &
  domainComponent in reverse order per RFC~4519 \\
\bottomrule
\end{tabularx}
\end{table*}

There are three combinators $\{\neg, \wedge, \vee\}$, with semantics following
classical propositional logic.

\textbf{Generality grading of atomic templates: GENERIC and NON\_GENERIC.} Atomic
templates are graded by two orthogonal criteria independently of the number of
parameters: an atomic template is GENERIC if and only if, first, it expresses a
general PKI concept---a class of certificate-attribute judgments reusable
across rules---and, second, it is parameterized over fields/values and bound to no
specific rule\_id or corpus-specific OID/single clause; otherwise it is
NON\_GENERIC. In that case, its logic is specialized to the internal structure of
some extension or a single RFC/CABF clause, and its semantics remain fixed to
that construct even when parameterized. A zero-parameter atomic template may
belong to either class. For example, IsCA is GENERIC, while
NotAfterIsNoExpirySentinel is NON\_GENERIC.
This grading is fixed in code as the two sets GENERIC\_ATOMS /
NON\_GENERIC\_ATOMS with a partition-completeness assertion, and is
deterministically recomputable.

Of the 151 atomic templates in $\mathcal{A}$, 98 are GENERIC and 53 are
NON\_GENERIC. Table~\ref{tab:nongeneric} lists only representative NON\_GENERIC
examples; the full list is released with the code and data.

\begin{table*}[tb]
\centering
\footnotesize
\caption{Representative NON\_GENERIC atomic templates.}
\label{tab:nongeneric}
\begin{tabularx}{\textwidth}{@{}p{4.6cm} p{3cm} X@{}}
\toprule
Representative NON\_GENERIC atomic template & Specialization target &
  Semantics \\
\midrule
SigAlgMatchesTBSSignature & certificate signature-algorithm fields &
  outer signatureAlgorithm byte-for-byte identical to
  tbsCertificate.signature \\
NotAfterIsNoExpirySentinel & validity sentinel value & whether
  notAfter is a specific no-expiry sentinel time \\
DomainComponentOrdered & Subject DN's domainComponent &
  whether the domainComponent sequence is ordered per RFC~4519 \\
CertPolicy\-ExplicitText\-HasEncodingTagInSet & certificatePolicies
  explicitText & whether explicitText uses an allowed ASN.1 string encoding \\
CRLDPHasNameRelative & CRL Distribution Points & whether
  distributionPoint uses the nameRelativeToCRLIssuer form \\
SubtreeIPListAny\-HasOctetCount\-AndNotAllZero & NameConstraints IP subtree
  & whether the IP subtree has an entry of given octet count that is not
  all-zero \\
WildcardFilter & DNS name patterns & whether a domain wildcard
  satisfies the specialized filter rule \\
\bottomrule
\end{tabularx}
\end{table*}

NON\_GENERIC atomic templates are allowed into code generation---they correspond
to real PKI constructs, not padding to inflate the corpus---but are explicitly
marked. The GENERIC/NON\_GENERIC split of the 93 generated lints is reported in
\S\ref{sec:eval}; among the 79 final-shipping unanimous lints the split is 60
GENERIC-only, 4 mixed, and 15 NON\_GENERIC-only.

\textbf{Rendering and metadata binding.} After $\rho$ renders them, all DSL trees
are embedded into a fixed zlint Go shell, with registration via
RegisterCertificateLint and the CheckApplies/Execute
method pair, so inter-rule differences are entirely concentrated in the check
body and imports; $\Phi_{\mathrm{post}}$ of \S\ref{subsec:synth} deterministically
binds Description/Citation/Name and each source's
PACKAGE/SOURCE/EFFECTIVE\_DATE metadata. For example,
RFC5280 $\mapsto$ RFC5280, CABF-TLS-BR $\mapsto$
CABFBaselineRequirements. The obligation-to-severity mapping is
MUST/MUST NOT/SHALL/SHALL NOT/REQUIRED $\mapsto$ lint.Error with lint-name
prefix e\_, and SHOULD/SHOULD NOT/RECOMMENDED $\mapsto$
lint.Warn with lint-name prefix w\_; MAY/OPTIONAL is already excluded by $C_1$ of
\S\ref{sec:lintability}, so the system produces no lint.Notice-level
output.

\section{Phrase Dictionary of the Mechanical Translation $\smech$}
\label{app:sigma}

$\smech$, defined in \S\ref{subsec:mech}, consists of an
atomic-template-to-phrase dictionary $\mathcal{M} : \mathcal{A} \to
\mathcal{L}_{\mathrm{NL}}$ and three combinator reduction rules. Representative
entries appear in Table~\ref{tab:phrasedict}.

\begin{table}[tb]
\centering
\footnotesize
\caption{Representative entries of the phrase dictionary $\mathcal{M}$.}
\label{tab:phrasedict}
\begin{tabularx}{\columnwidth}{@{}p{3.1cm} X@{}}
\toprule
Atomic template & $\mathcal{M}(a)$ phrase template \\
\midrule
ExtPresent(O) & ``the \{O\} extension is present'' \\
ExtCritical(O) & ``the \{O\} extension is present and marked
  critical'' \\
FieldEq(F, v) & ``\{F\} equals \{v\}'' \\
FieldEncodedAs(F, T) & ``\{F\} is encoded as ASN.1 \{T\}'' \\
KeyUsageHas(B) & ``the KeyUsage bit \{B\} is asserted'' \\
ListAllMatch(F, T) & ``every entry of \{F\} satisfies
  the result of $\smech(T)$'' \\
SubtreeIPListAny\-HasOctetCount\-AndNotAllZero(F, n) & ``the
  NameConstraints IP subtree \{F\} contains a non-zero entry of \{n\}
  octets'' \\
\bottomrule
\end{tabularx}
\end{table}

Combinator reduction: $\smech(\neg t) = $ ``NOT ($\smech(t)$)'';
$\smech(t_1 \wedge t_2) = $ ``($\smech(t_1)$) AND ($\smech(t_2)$)''; $\vee$
likewise; the atomic-template base case $\smech(a) = \mathcal{M}(a)$. For a
conditional pair $(p,q)$: when $p = \neg(p')$, it outputs ``WHEN NOT
($\smech(p')$), THEN $\smech(q)$'', the NEG-PRE template that eliminates double
negation; otherwise it outputs ``WHEN ($\smech(p)$), THEN $\smech(q)$''; and
when $p = \varepsilon$ it outputs the single sentence $\smech(q)$.

%-------------------------------------------------------------------------------
% Stage-Aware Iterative Verification, abbreviated SAIV --- moved out of the main body and
% merged inline here; previously the separate file 07_saiv.tex, now deleted.
%-------------------------------------------------------------------------------
\section{Stage-Aware Iterative Verification: SAIV}
\label{sec:saiv}

To coordinate the modules and safeguard their output quality without manual
ground truth, we propose the Stage-Aware Iterative Verification framework,
abbreviated SAIV.
The basic idea is to predefine two quantifiable optimization
objectives, measure the gap between the current and ideal states by each
objective's computable residual, and when a residual has not reached its
threshold, the operator uses the residual signal to locate the offending module
and applies the corresponding repair, then re-measures and proceeds to the next
iteration, until both residuals converge. Alongside the two objectives, two
auxiliary checks police the lintability decision without entering this loop.

\subsection{Formal Definition}
\label{subsec:saiv-formal}

The whole pipeline $\Pi$ chains four modules---extraction $\to$ classification
$\to$ generation $\to$ verification:
\begin{equation}\label{eq:pipeline}
  \Pi = \phi_V \circ \phi_G \circ \phi_C \circ \phi_R
\end{equation}
where $\phi_R$ is extraction, $\phi_C$ is the four-condition lintability decision
of \S\ref{sec:lintability}, $\phi_G$ is the DSL code space $\Tv$ of
\S\ref{sec:codegen}, and $\phi_V$ is the three-layer semantic alignment of
\S\ref{sec:codegen}. $\phi_C$ is the
single-rule lintability predicate $r \mapsto \{0,1\}$; the $\psi_C$ of
Algorithm~\ref{alg:saiv} is its set-level lift, which strips the noise $\RN$ from
$\Rkw$ and sorts the remaining genuine rules into lintable $\RL$ and not-lintable
$\RU$, giving the partition $(\RN, \RL, \RU)$. The operator drives the loop from
the computable signals these modules produce.

\subsection{Two Iteration Objectives and Two Auxiliary Lintability Checks}
\label{subsec:targets}

SAIV borrows the idea of backpropagation. It fixes two optimization objectives
the whole pipeline is trained to approach, and keeps apart from them two
auxiliary checks that police the lintability decision $\phi_C$ rather than drive
optimization. The two objectives G1 and G2 each carry a directly computable
residual, $\Lrec$ and $\Lcode$, whose weighted sum is the total loss $\Ltot$, so
convergence is defined by $\Ltot$ alone. S1 and S2 are checks, not objectives,
and add no term to $\Ltot$. They form the two sides of a single consistency test
on $\phi_C$, S1 auditing the not-lintable verdict and S2 the lintable verdict,
each firing only to correct a suspected misclassification. The subsequent subsections
revolve around the two objectives, while the two checks gate the labels those
objectives consume.

\medskip
\noindent\textbf{G1: recall completeness.}
\begin{theorem}[Recall-Completeness Invariant]\label{thm:recall}
If the classification $\psi_C$ assigns each keyword-recalled candidate in
$\Rkw(\mathcal{D})$ to exactly one of the three classes---noise, lintable, and
not-lintable---then recall is conserved from recall to classification:
\begin{equation}\label{eq:cons}
  |\Rkw(\mathcal{D})| = |\RN| + |\RL| + |\RU|
\end{equation}
\end{theorem}
\noindent where $\mathcal{D}$ is the standard-document corpus, written D in
ASCII inside the Algorithm~\ref{alg:saiv} listing, $\RN$ is noise that contains
RFC~2119 keywords but not genuine rules, $\RL$ the lintable rules, and $\RU$
the not-lintable rules. The executable class $\RL$ is further split by ``whether
it is already implemented by an existing external lint tool'':
\begin{equation}\label{eq:split}
  |\RL| = |\RLcov| + |\RLunc|
\end{equation}
where $\RLcov$ are rules covered by some lint tool and $\RLunc$ are
not-yet-covered rules. The first identity~\eqref{eq:cons} is total conservation
from recall to classification, and the second~\eqref{eq:split} splits the
executable set so the later objectives can be pursued separately. The uncovered
subset $\RLunc$ is the most important observation domain for generation that
``fills the ecosystem gap,'' while the covered subset $\RLcov$ provides external
evidence of reverse executability. The residual is the recall-conservation
residual $\Lrec$, a symmetric normalization of the gap in
Theorem~\ref{thm:recall}, computable without manual ground truth and currently
closed, as shown in Table~\ref{tab:snapshot}:
\begin{equation}\label{eq:recall-res}
  \Lrec = 1 - \frac{\min\!\big(|\Rkw|,\; |\RN|+|\RL|+|\RU|\big)}
  {\max\!\big(|\Rkw|,\; |\RN|+|\RL|+|\RU|\big)}
\end{equation}

\medskip
\noindent\textbf{G2: code--specification synonymy for uncovered rules.} For the
uncovered subset $\RLunc$, the code summary $\smech(t_r)$, the code\_summary of
rule $r$'s DSL tree $t_r$ computed by $\smech$ of \S\ref{subsec:mech}, is
required to be synonymous with the specification text.
The residual $\Lcode$ is the average, over $\RLunc$, of the per-rule correctness
label $\lambda_{\mathrm{code}}(r) = \ind[\mathrm{compile}(\rho(t_r))] \cdot
\mathrm{syn}(\smech(t_r), \mathrm{spec}(r))$, where $\rho(t_r)$ is the rendered Go
code and $\mathrm{spec}(r)$ the rule text. The structural-completeness score is
$\equiv 1$ in the fixed zlint shell, so it drops out:
\begin{equation}\label{eq:code-res}
  \Lcode = 1 - \frac{1}{|\RLunc|}\sum_{r\in\RLunc}\lambda_{\mathrm{code}}(r)
\end{equation}

\medskip
\noindent\textbf{S1: reverse executability of covered rules.} The first
auxiliary check, auditing the not-lintable verdict. If a rule has a corresponding
implementation in a lint tool such as zlint it must be lintable, so the check
counts the violations $\Nviol$ and is clear when $\Nviol = 0$. The operator
computes $\Nviol$ by a reverse coverage check---running the coverage judge over
the not-lintable set $\RU$ and counting the rules it rates full or partial; each
such hit is adjudicated as either a $\phi_C$ false negative and re-extracted to
lintable, or a spurious coverage match and kept not-lintable.

\medskip
\noindent\textbf{S2: synthesizability of lintable rules.} The dual check,
auditing the lintable verdict. A rule judged lintable should be synthesizable
within $\Tv$, so a no-template abstention $\bot_{\mathrm{NT}}$ returned for such a
rule after generation repair $\mathcal{P}_G$ and vocabulary expansion
$\mathcal{P}_A$ are exhausted is a candidate signal that $\phi_C$ mislabeled it.
S2 adds no residual to $\Ltot$; it reuses the generator's existing abstention
marker as a diagnostic, and each persistent abstention is reviewed against the
four-condition decision of \S\ref{sec:lintability} and adjudicated as either a
$\phi_C$ false positive and re-derived to not-lintable, or a genuine $\Tv$
expressiveness gap and kept as a logged abstention. In the shipping corpus both
no-template abstentions are $\Tv$ gaps, so neither is re-adjudicated as
not-lintable.

\medskip
The two objective residuals are weighted into the total loss in
Eq.~\eqref{eq:total}, with $w_R + w_C = 1$:
\begin{equation}\label{eq:total}
  \Ltot = w_R\cdot\Lrec + w_C\cdot\Lcode
\end{equation}
Since recall is conserved, $\Lrec = 0$, $\Ltot$ is governed by $\Lcode$ alone.
The two auxiliary checks stay outside $\Ltot$; their signals are the violation
count $\Nviol$ of S1 and the abstention review of S2 in \S\ref{subsec:algo}.

\begin{table*}[tb]
\centering
\footnotesize
\caption{The two SAIV objectives and the two auxiliary lintability checks, with
their signals and repair routing.}
\label{tab:saiv-objectives}
\begin{tabularx}{\textwidth}{@{}llp{4.0cm}X@{}}
\toprule
Item & Signal & Main repair direction & Repair target \\
\midrule
\multicolumn{4}{@{}l}{\emph{Iteration objectives --- define $\Ltot$ and
  convergence}}\\
G1 Recall completeness & $\Lrec$, Eq.~\eqref{eq:recall-res} & --- &
  recall logic \\
G2 Synonymy & $\Lcode$, Eq.~\eqref{eq:code-res} &
  IR-content, DSL-generation, and verification-judgment repairs &
  IR-content correction / DSL-tree re-synthesis / judgment \& summary
  re-check \\
\midrule
\multicolumn{4}{@{}l}{\emph{Auxiliary lintability checks --- police $\phi_C$,
  outside $\Ltot$}}\\
S1 Reverse executability & $\Nviol$, \S\ref{subsec:algo} &
  verification-judgment consistency repair &
  $\phi_C$ false negatives / spurious coverage matches \\
S2 Synthesizability & $\bot_{\mathrm{NT}}$ abstention, \S\ref{subsec:algo} &
  generation repair then vocabulary expansion &
  $\phi_C$ false positives / $\Tv$ expressiveness gaps \\
\bottomrule
\end{tabularx}
\end{table*}

Table~\ref{tab:saiv-objectives} gives only the item-to-repair-direction correspondence; the
concrete repair operations and their notation are defined in
\S\ref{subsec:repair}.

\subsection{Stage Routing and Repair Operations}
\label{subsec:repair}

\textbf{Stage routing.} Each round the operator chooses the next repair operation using
directly computable quantities, namely the recall-conservation residual, the
compilation-failure rate over the lintable set, the average structural score,
the entity-necessary-condition screening result, and the synonymy confidence.
This routing only decides which kind of repair to attempt next round.

\textbf{Repair operations.} SAIV comprises four kinds of operations---IR-content
repair $\mathcal{P}_R$, classification repair $\mathcal{P}_C$, generation repair
$\mathcal{P}_G$, and verification repair $\mathcal{P}_V$, the subscripts naming
the repaired module $\phi_\bullet$. The repair operator $\mathcal{P}$ is distinct
from the rendering function $\rho$ of \S\ref{subsec:mech}. The operator selects
the next round's operation from the signals above, and a repair is accepted only
when the metric decreases.

\begin{itemize}
\item \textbf{IR-content repair $\mathcal{P}_R$}: feeds the failure trace---the
  current IR, the irreducible category, the mechanical summary, the judge ruling
  and rationale, the synonymy confidence, and the session's IR history---back to
  the LLM, which returns REPAIR(ir') or NO\_FIX.
\item \textbf{Classification repair $\mathcal{P}_C$}: re-extracts the
  four-condition IR fields and re-runs the lintability decision.
\item \textbf{Generation repair $\mathcal{P}_G$}: tries deterministic repair
  first; on failure it feeds the failing code fragment and the IR-constraint
  discrepancy, or the parser's signature/closure error, back into the prompt to
  re-synthesize within $\Tv$; on persistent abstention it enters the offline
  vocabulary-expansion channel $\mathcal{P}_A$, which hand-authors a new
  atomic-template candidate and merges it into $\mathcal{A}$ after
  certificate-level execution verification.
\item \textbf{Verification repair $\mathcal{P}_V$}: enlarges the synonymy
  judgment's semantic neighborhood, e.g., negation/affirmation interchange or
  one-to-many decomposition, or corrects the inconsistency between the two
  verdict sources.
\end{itemize}

\subsection{Iteration Algorithm and Termination Conditions}
\label{subsec:algo}

Algorithm~\ref{alg:saiv} gives the full procedure.

% Algorithm 2 -- Stage-Aware Iterative Verification, abbreviated SAIV.
% Included from the SAIV section of appendix.tex via % Algorithm 2 -- Stage-Aware Iterative Verification, abbreviated SAIV.
% Included from the SAIV section of appendix.tex via % Algorithm 2 -- Stage-Aware Iterative Verification, abbreviated SAIV.
% Included from the SAIV section of appendix.tex via \input{algorithms/alg_saiv}.
%
% Typeset single-column (not figure*) at \scriptsize, for the same reason as
% Algorithm 1: short algorithmic lines wasted the right half of a full-width float,
% which is billed two columns of page area. The RouteRepair \Comment is on its own
% \Statex line because \Comment is right-aligned and cannot wrap. No step, symbol,
% or comment wording was changed.
\begin{figure}[tb]
\setlength{\parskip}{0pt}
\makeatletter
\@ifundefined{c@algorithm}{%
  \newcounter{algorithm}%
  \renewcommand{\thealgorithm}{\arabic{algorithm}}%
}{}%
\makeatother
\refstepcounter{algorithm}
\label{alg:saiv}
\hrule width \linewidth height 0.6pt\relax
\vspace{1pt}
\nointerlineskip
\noindent\textbf{Algorithm~\thealgorithm.} Algorithmic listing for
Stage-Aware Iterative Verification, abbreviated SAIV.
\par
\vspace{1pt}
\nointerlineskip
\hrule width \linewidth height 0.4pt\relax
\nointerlineskip
\scriptsize
\begin{algorithmic}[1]
\Statex \textbf{Input:} Web PKI normative-source corpus $D$, threshold $\theta$, max iterations $K$
\Statex \textbf{Output:} final code set $C^{*}$ and termination status
\State $\Rkw \gets \phi_R(D)$;\quad $(\RN,\RL,\RU)\gets \psi_L(\Rkw)$
  \Comment{$\RU=$ not-lintable}
\State $(\RLcov,\RLunc)\gets \textsc{Coverage}(\RL)$
  \Comment{Algorithm~\ref{alg:coverage}}
\State $C \gets \{\phi_G(r): r\in \RLunc\}$
\State $t\gets 0$
\Repeat
  \State compute objective residuals $\Lrec$~\eqref{eq:recall-res},
    $\Lcode$~\eqref{eq:code-res}; lintability-check signals $\Nviol$ (S1) and the
    pending-abstention set $\NTpend$ (S2)
  \If{$\Ltot<\theta$ \textbf{and} $\Nviol=0$ \textbf{and} $|\NTpend|=0$}
    \State \textbf{break} \Comment{objectives converged, S1 and S2 clear}
  \EndIf
  \State $\mathit{op}\gets \textsc{RouteRepair}(\Lrec,\Lcode,\Nviol,\NTpend,\cdots)$
  \Statex \Comment{repair routing, \S\ref{subsec:repair}}
  \If{$\mathit{op}=\mathcal{P}_R$} \Comment{IR-content repair}
    \ForAll{$r$ routed to $\mathcal{P}_R$}
      \State $\mathit{rep}\gets \mathcal{P}_R(\textsc{FailTrace}(r))$
        \Comment{REPAIR(ir$'$) or NO\_FIX}
      \If{$\mathit{rep}=\textsc{No\_Fix}$ \textbf{or}
          \textsc{ReextractCheck}$(\mathit{rep}.ir',\textsc{text}(r))\neq\emptyset$}
        \State record $r$ in irreducible set $\mathcal{R}_{\mathrm{irred}}$
          \Comment{kept as residual}
      \ElsIf{\textsc{Retest}$(\mathit{rep}.ir')$ passes}
        \State $\mathrm{IR}(r)\gets \mathit{rep}.ir'$ \Comment{accept}
      \EndIf
    \EndFor
  \Else
    \State apply $\mathit{op}\in\{\mathcal{P}_L,\mathcal{P}_G,\mathcal{P}_V\}$
  \EndIf
  \State update $(\RN,\RL,\RU,\RLunc)$ and $C$
  \State $t\gets t+1$
\Until{$t\ge K$ \textbf{or} no residual decreased, the no-progress condition}
\State \textbf{return} $(C,\ \mathit{status})$, with $\mathit{status}=\textsc{closed}$
  on break, else \textsc{max\_iter} if $t\ge K$ and \textsc{dry} otherwise
\end{algorithmic}
\vspace{1pt}
\hrule width \linewidth height 0.6pt\relax
\end{figure}
.
%
% Typeset single-column (not figure*) at \scriptsize, for the same reason as
% Algorithm 1: short algorithmic lines wasted the right half of a full-width float,
% which is billed two columns of page area. The RouteRepair \Comment is on its own
% \Statex line because \Comment is right-aligned and cannot wrap. No step, symbol,
% or comment wording was changed.
\begin{figure}[tb]
\setlength{\parskip}{0pt}
\makeatletter
\@ifundefined{c@algorithm}{%
  \newcounter{algorithm}%
  \renewcommand{\thealgorithm}{\arabic{algorithm}}%
}{}%
\makeatother
\refstepcounter{algorithm}
\label{alg:saiv}
\hrule width \linewidth height 0.6pt\relax
\vspace{1pt}
\nointerlineskip
\noindent\textbf{Algorithm~\thealgorithm.} Algorithmic listing for
Stage-Aware Iterative Verification, abbreviated SAIV.
\par
\vspace{1pt}
\nointerlineskip
\hrule width \linewidth height 0.4pt\relax
\nointerlineskip
\scriptsize
\begin{algorithmic}[1]
\Statex \textbf{Input:} Web PKI normative-source corpus $D$, threshold $\theta$, max iterations $K$
\Statex \textbf{Output:} final code set $C^{*}$ and termination status
\State $\Rkw \gets \phi_R(D)$;\quad $(\RN,\RL,\RU)\gets \psi_L(\Rkw)$
  \Comment{$\RU=$ not-lintable}
\State $(\RLcov,\RLunc)\gets \textsc{Coverage}(\RL)$
  \Comment{Algorithm~\ref{alg:coverage}}
\State $C \gets \{\phi_G(r): r\in \RLunc\}$
\State $t\gets 0$
\Repeat
  \State compute objective residuals $\Lrec$~\eqref{eq:recall-res},
    $\Lcode$~\eqref{eq:code-res}; lintability-check signals $\Nviol$ (S1) and the
    pending-abstention set $\NTpend$ (S2)
  \If{$\Ltot<\theta$ \textbf{and} $\Nviol=0$ \textbf{and} $|\NTpend|=0$}
    \State \textbf{break} \Comment{objectives converged, S1 and S2 clear}
  \EndIf
  \State $\mathit{op}\gets \textsc{RouteRepair}(\Lrec,\Lcode,\Nviol,\NTpend,\cdots)$
  \Statex \Comment{repair routing, \S\ref{subsec:repair}}
  \If{$\mathit{op}=\mathcal{P}_R$} \Comment{IR-content repair}
    \ForAll{$r$ routed to $\mathcal{P}_R$}
      \State $\mathit{rep}\gets \mathcal{P}_R(\textsc{FailTrace}(r))$
        \Comment{REPAIR(ir$'$) or NO\_FIX}
      \If{$\mathit{rep}=\textsc{No\_Fix}$ \textbf{or}
          \textsc{ReextractCheck}$(\mathit{rep}.ir',\textsc{text}(r))\neq\emptyset$}
        \State record $r$ in irreducible set $\mathcal{R}_{\mathrm{irred}}$
          \Comment{kept as residual}
      \ElsIf{\textsc{Retest}$(\mathit{rep}.ir')$ passes}
        \State $\mathrm{IR}(r)\gets \mathit{rep}.ir'$ \Comment{accept}
      \EndIf
    \EndFor
  \Else
    \State apply $\mathit{op}\in\{\mathcal{P}_L,\mathcal{P}_G,\mathcal{P}_V\}$
  \EndIf
  \State update $(\RN,\RL,\RU,\RLunc)$ and $C$
  \State $t\gets t+1$
\Until{$t\ge K$ \textbf{or} no residual decreased, the no-progress condition}
\State \textbf{return} $(C,\ \mathit{status})$, with $\mathit{status}=\textsc{closed}$
  on break, else \textsc{max\_iter} if $t\ge K$ and \textsc{dry} otherwise
\end{algorithmic}
\vspace{1pt}
\hrule width \linewidth height 0.6pt\relax
\end{figure}
.
%
% Typeset single-column (not figure*) at \scriptsize, for the same reason as
% Algorithm 1: short algorithmic lines wasted the right half of a full-width float,
% which is billed two columns of page area. The RouteRepair \Comment is on its own
% \Statex line because \Comment is right-aligned and cannot wrap. No step, symbol,
% or comment wording was changed.
\begin{figure}[tb]
\setlength{\parskip}{0pt}
\makeatletter
\@ifundefined{c@algorithm}{%
  \newcounter{algorithm}%
  \renewcommand{\thealgorithm}{\arabic{algorithm}}%
}{}%
\makeatother
\refstepcounter{algorithm}
\label{alg:saiv}
\hrule width \linewidth height 0.6pt\relax
\vspace{1pt}
\nointerlineskip
\noindent\textbf{Algorithm~\thealgorithm.} Algorithmic listing for
Stage-Aware Iterative Verification, abbreviated SAIV.
\par
\vspace{1pt}
\nointerlineskip
\hrule width \linewidth height 0.4pt\relax
\nointerlineskip
\scriptsize
\begin{algorithmic}[1]
\Statex \textbf{Input:} Web PKI normative-source corpus $D$, threshold $\theta$, max iterations $K$
\Statex \textbf{Output:} final code set $C^{*}$ and termination status
\State $\Rkw \gets \phi_R(D)$;\quad $(\RN,\RL,\RU)\gets \psi_L(\Rkw)$
  \Comment{$\RU=$ not-lintable}
\State $(\RLcov,\RLunc)\gets \textsc{Coverage}(\RL)$
  \Comment{Algorithm~\ref{alg:coverage}}
\State $C \gets \{\phi_G(r): r\in \RLunc\}$
\State $t\gets 0$
\Repeat
  \State compute objective residuals $\Lrec$~\eqref{eq:recall-res},
    $\Lcode$~\eqref{eq:code-res}; lintability-check signals $\Nviol$ (S1) and the
    pending-abstention set $\NTpend$ (S2)
  \If{$\Ltot<\theta$ \textbf{and} $\Nviol=0$ \textbf{and} $|\NTpend|=0$}
    \State \textbf{break} \Comment{objectives converged, S1 and S2 clear}
  \EndIf
  \State $\mathit{op}\gets \textsc{RouteRepair}(\Lrec,\Lcode,\Nviol,\NTpend,\cdots)$
  \Statex \Comment{repair routing, \S\ref{subsec:repair}}
  \If{$\mathit{op}=\mathcal{P}_R$} \Comment{IR-content repair}
    \ForAll{$r$ routed to $\mathcal{P}_R$}
      \State $\mathit{rep}\gets \mathcal{P}_R(\textsc{FailTrace}(r))$
        \Comment{REPAIR(ir$'$) or NO\_FIX}
      \If{$\mathit{rep}=\textsc{No\_Fix}$ \textbf{or}
          \textsc{ReextractCheck}$(\mathit{rep}.ir',\textsc{text}(r))\neq\emptyset$}
        \State record $r$ in irreducible set $\mathcal{R}_{\mathrm{irred}}$
          \Comment{kept as residual}
      \ElsIf{\textsc{Retest}$(\mathit{rep}.ir')$ passes}
        \State $\mathrm{IR}(r)\gets \mathit{rep}.ir'$ \Comment{accept}
      \EndIf
    \EndFor
  \Else
    \State apply $\mathit{op}\in\{\mathcal{P}_L,\mathcal{P}_G,\mathcal{P}_V\}$
  \EndIf
  \State update $(\RN,\RL,\RU,\RLunc)$ and $C$
  \State $t\gets t+1$
\Until{$t\ge K$ \textbf{or} no residual decreased, the no-progress condition}
\State \textbf{return} $(C,\ \mathit{status})$, with $\mathit{status}=\textsc{closed}$
  on break, else \textsc{max\_iter} if $t\ge K$ and \textsc{dry} otherwise
\end{algorithmic}
\vspace{1pt}
\hrule width \linewidth height 0.6pt\relax
\end{figure}


\textbf{Termination conditions}: $\Ltot < \theta$, with default $0.05$, and the two
lintability checks clear, reaching the maximum iteration $K = 10$, or no residual
decreasing in a round. If every
accepted stage-repair operation is monotone, reducing that stage's error or
leaving it unchanged, then $\Ltot$ is non-increasing each round; the actual
system uses a ``stop when no decrease'' policy, so it guarantees finite-round
termination but not global optimality, as discussed in \S\ref{sec:eval}.
